\chapter{System Design and Architecture}
\label{ch:architecture}

This chapter presents the design of \zkwns. It begins with the design goals and
the threat model that constrain every subsequent decision, then gives a layered
overview of the architecture and the rationale for its defining choice---moving
key custody from hardware to mathematics. It then specifies each of the three
functional modules, the data model (with particular attention to what is and is
not placed on-chain), and the end-to-end workflows that tie the modules together.
Implementation details are deferred to Chapter~\ref{ch:implementation}; this
chapter is concerned with \emph{what} the system does and \emph{why}.

\section{Design Goals and Principles}
\label{sec:design-goals}

The architecture is shaped by six principles, each traceable to the problem
statement of Section~\ref{sec:intro-problem}.

\begin{description}
  \item[P1 --- No single point of trust.] No individual party---person,
  organisation, server, or chip---may be in a position to unilaterally read the
  evidence, withhold it, release it early, or deanonymise the source. Every
  trust assumption is to be distributed across an independent network.

  \item[P2 --- Client-side confidentiality.] Plaintext and key material exist
  only on the source's device. Everything that leaves the browser is either
  ciphertext, a commitment, or a public reference.

  \item[P3 --- No plaintext on-chain.] The public ledger stores only the minimum
  needed to coordinate the protocol: timestamps, references to off-chain data,
  hashes, and payment state. It never stores sensitive content.

  \item[P4 --- Autonomy and unstoppability.] The release of evidence must not
  depend on any operator's continued cooperation. Once armed, the switch fires
  based solely on objective, on-chain state.

  \item[P5 --- Credible anonymity.] The source can substantiate a claim about
  themselves (e.g.\ insider access) to the level a journalist needs, while
  remaining anonymous; reward must not break that anonymity.

  \item[P6 --- Graceful degradation and honesty.] Where a component cannot
  provide its full guarantee in a given environment, the system must say so
  clearly rather than fail silently or pretend.
\end{description}

\section{Threat Model}
\label{sec:design-threat}

\paragraph{Actors.} The protagonists are the \emph{source} (whistleblower), who
holds evidence and wishes to seal it; the \emph{recipient} (a specific
journalist) or the \emph{public}, to whom the evidence is released; and a
\emph{buyer} (journalist) in the marketplace, who pays for access.

\paragraph{Adversary.} The adversary is powerful: a well-resourced organisation,
up to and including a nation-state, that may (i) surveil network traffic and the
public blockchain; (ii) compel or compromise a centralised service provider, a
cloud host, or a hardware vendor under legal or extralegal pressure; (iii) detain
or coerce the source; and (iv) attempt to submit fabricated leaks or to
deanonymise sources through metadata. The adversary is assumed \emph{not} able to
break standard cryptographic primitives (AES-256, secp256k1, BLS, SHA-256) or to
simultaneously corrupt a threshold fraction of an independent, globally
distributed node network.

\paragraph{Security goals.} Against this adversary the system aims to guarantee:
confidentiality of the evidence until the release condition holds; integrity and
authenticity of the released evidence; autonomy of the release (no operator can
prevent it once armed); anonymity of the source throughout; unlinkability of
rewards; and a credibility signal that resists Sybil and fabrication attacks.

\paragraph{Explicitly out of scope.} The system does not defend against
compromise of the source's own device while they hold plaintext (an endpoint
problem), against the source voluntarily revealing their identity, against
traffic-analysis deanonymisation at the network layer (for which the source
should use Tor or a VPN), or against the legal consequences a source may face if
independently identified. These are acknowledged in Chapter~\ref{ch:evaluation}.

\section{Architectural Overview}
\label{sec:design-overview}

\zkw is structured as four cooperating layers, each operated by a different
decentralised system so that no layer can compromise the guarantees of another.
Figure~\ref{fig:arch-layers} shows the layers and the data that flows between
them.

\begin{figure}[htbp]
\centering
\begin{tikzpicture}[
  font=\sffamily\small,
  layer/.style={draw=zkwframe, rounded corners, fill=zkwbg, minimum width=12.2cm, minimum height=1.55cm, align=center},
  box/.style={draw=zkwaccent, rounded corners, fill=white, align=center, minimum height=0.95cm, font=\sffamily\scriptsize, inner sep=3pt},
  lbl/.style={font=\sffamily\bfseries\footnotesize, text=zkwaccent},
  >={Stealth[length=2mm]}
]
% Client layer
\node[layer, minimum height=1.8cm] (client) {};
\node[lbl, anchor=north west] at ([xshift=4pt,yshift=-3pt]client.north west) {Client Layer (in-browser, Next.js)};
\node[box] (enc) at ([xshift=-3.9cm,yshift=-0.25cm]client.center) {AES-256-GCM\\encryption};
\node[box] (rec) at ([xshift=0cm,yshift=-0.25cm]client.center) {Reclaim\\zkTLS widget};
\node[box] (ste) at ([xshift=3.9cm,yshift=-0.25cm]client.center) {Stealth keys\\(secp256k1)};

% Chain layer
\node[layer, below=1.3cm of client, minimum height=1.8cm] (chain) {};
\node[lbl, anchor=north west] at ([xshift=4pt,yshift=-3pt]chain.north west) {Logic / State Layer (EVM smart contracts)};
\node[box] (dms) at ([xshift=-3.9cm,yshift=-0.25cm]chain.center) {DeadMansSwitch\\isDeceased()};
\node[box] (reg) at ([xshift=0cm,yshift=-0.25cm]chain.center) {Whistleblower\\Registry};
\node[box] (mkt) at ([xshift=3.9cm,yshift=-0.25cm]chain.center) {Marketplace\\(stealth payout)};

% Keys layer
\node[layer, below=1.3cm of chain] (keys) {};
\node[lbl, anchor=north west] at ([xshift=4pt,yshift=-3pt]keys.north west) {Key-Management Layer};
\node[box, minimum width=6cm] (lit) at ([yshift=-0.2cm]keys.center) {Lit Protocol MPC / threshold-BLS network\\(Access Control Conditions)};

% Storage layer
\node[layer, below=1.3cm of keys] (store) {};
\node[lbl, anchor=north west] at ([xshift=4pt,yshift=-3pt]store.north west) {Storage Layer};
\node[box, minimum width=6cm] (arw) at ([yshift=-0.2cm]store.center) {Arweave permaweb\\(encrypted payload, permanent)};

% Arrows
\draw[->, zkwaccent, thick] (enc.south) -- ++(0,-0.55) -| (arw.north) node[midway, above, pos=0.25, font=\sffamily\tiny\itshape, text=zkwcomment] {ciphertext};
\draw[->, zkwaccent, thick] (enc.south) -- ++(0,-0.3) -| (lit.north) node[near start, above, font=\sffamily\tiny\itshape, text=zkwcomment] {AES key};
\draw[->, dashed, zkwaccent] (lit.east) -- ++(1.0,0) |- (dms.east) node[near start, right, font=\sffamily\tiny\itshape, text=zkwcomment] {reads};
\draw[->, zkwaccent, thick] (rec.south) -- (reg.north) node[midway, right, font=\sffamily\tiny\itshape, text=zkwcomment] {proof hash};
\draw[->, zkwaccent, thick] (ste.south) -- (mkt.north) node[midway, right, font=\sffamily\tiny\itshape, text=zkwcomment] {meta-addr.};
\end{tikzpicture}
\caption{The four-layer architecture of \zkwns. Each layer is operated by an
independent decentralised system; arrows show the principal data flows.}
\label{fig:arch-layers}
\end{figure}

The \textbf{client layer} is a Next.js single-page application that runs entirely
in the source's browser. It performs all encryption, drives the zkTLS proof
widget, and manages stealth-address keys. By principle P2, it is the only place
plaintext ever exists.

The \textbf{logic/state layer} comprises three EVM smart contracts. They hold no
sensitive data; they coordinate the protocol by recording liveness
(\texttt{DeadMansSwitch}), reputation commitments
(\texttt{WhistleblowerRegistry}), and marketplace state (\texttt{Marketplace}).

The \textbf{key-management layer} is the Lit Protocol threshold network. It seals
the AES key and releases it only when the dead man's switch's on-chain condition
evaluates true. It is the architectural replacement for the V1 TEE.

The \textbf{storage layer} is Arweave, which holds the encrypted payload
permanently and censorship-resistantly.

The separation matters because the layers have \emph{different and independent
failure modes}. Compromising the chain does not reveal the key; compromising
storage yields only ciphertext; compromising a minority of Lit nodes reveals
nothing; and none of them ever sees plaintext. An adversary must defeat several
independent systems at once---the concrete realisation of principle P1.

\section{From Hardware to Mathematics}
\label{sec:design-shift}

The single most consequential design decision is the replacement of the V1
trusted execution environment with the Lit Protocol's threshold network as the
custodian of decryption keys. Section~\ref{sec:bg-tee} laid out the cryptographic
argument; here it is framed as an architectural one.

A TEE and a threshold network both answer the question ``who can hold the key
without being able to abuse it?''---but they answer it with different kinds of
assumption. The TEE answers with a \emph{physical} assumption: that a particular
piece of silicon cannot be opened, that its manufacturer is honest and
uncompelled, and that the specific nodes hosting the enclave stay online. The
threshold network answers with a \emph{distributional} assumption: that an
adversary cannot simultaneously corrupt a threshold fraction of many independent
operators. The crucial difference is what happens when the assumption is stressed
by a nation-state adversary. A physical assumption fails \emph{all at once and
silently}: one successful side-channel attack, or one compelled vendor, breaks
every key. A distributional assumption \emph{degrades gracefully and visibly}:
corrupting one operator, or ten, achieves nothing until the threshold is reached,
and doing so requires coordinated action across jurisdictions.

In the architecture, this manifests concretely. The AES key is encrypted such
that the Lit network will only assemble the decryption shares when an
\emph{Access Control Condition} holds. The condition is not a trusted operator's
say-so but a read of public, objective on-chain state:
\texttt{DeadMansSwitch.isDeceased(source)} must return \texttt{true}. The release
logic is thus split between an immutable smart contract (the condition) and a
threshold network (the key), neither of which can act alone. This is the design's
answer to principle P4.

Chapter~\ref{ch:evaluation} revisits this choice in light of the fact that the
Lit ecosystem has itself begun to offer a TEE-based runtime; the design here
deliberately uses the multi-party-computation network generation precisely to
keep the trust anchor in mathematics.

\section{Module A: The Vault (Dead Man's Switch)}
\label{sec:design-vault}

Module A is the heart of the system. Its on-chain component, the
\texttt{DeadMansSwitch} contract, is a minimal liveness registry. Each source has
at most one active \emph{switch} recording: the timestamp of their last check-in,
the required interval between check-ins, a reference to the encrypted payload on
Arweave, a reference to the Lit access-control conditions, and an optional
recipient address (the zero address denoting public release).

The switch exposes the predicate that the whole system pivots on:

\[
\texttt{isDeceased}(u) \;=\;
\begin{cases}
\texttt{true}  & \text{if the switch is active and } t_{\text{now}} > t_{\text{last}} + \Delta,\\[2pt]
\texttt{false} & \text{otherwise,}
\end{cases}
\]

where $t_{\text{last}}$ is the last check-in, $\Delta$ the interval, and
$t_{\text{now}}$ the current block timestamp. The source keeps the switch
``alive'' by periodically calling \texttt{checkIn()}, which resets $t_{\text{last}}$
to the current time. They may \texttt{deactivateSwitch()} to cancel while alive,
or \texttt{updateSwitchMetadata()} to re-point the storage and key references.

\begin{figure}[htbp]
\centering
\begin{tikzpicture}[
  font=\sffamily\small,
  node distance=2.4cm and 3.0cm,
  state/.style={draw=zkwaccent, thick, rounded corners, fill=zkwbg, minimum width=2.6cm, minimum height=1.1cm, align=center},
  term/.style={draw=zkwcomment, thick, rounded corners, fill=white, minimum width=2.6cm, minimum height=1.1cm, align=center},
  >={Stealth[length=2.2mm]}
]
\node[state] (none) {No switch};
\node[state, right=of none] (active) {Active\\(alive)};
\node[term, right=of active] (triggered) {Triggered\\(released)};
\node[term, below=1.6cm of active] (inactive) {Deactivated};

\draw[->, thick] (none) -- node[above, font=\sffamily\scriptsize] {createSwitch} (active);
\draw[->, thick] (active) to[loop above] node[above, font=\sffamily\scriptsize] {checkIn (resets timer)} (active);
\draw[->, thick] (active) -- node[above, font=\sffamily\scriptsize, align=center] {interval elapses\\(isDeceased = true)} (triggered);
\draw[->, thick] (active) -- node[right, font=\sffamily\scriptsize] {deactivateSwitch} (inactive);
\draw[->, thick, dashed] (inactive) -- node[below left, font=\sffamily\scriptsize, align=center] {createSwitch\\(re-arm)} (active);
\end{tikzpicture}
\caption{Lifecycle of a Dead Man's Switch. The transition to \emph{Triggered}
happens automatically with the passage of time; it requires no transaction.}
\label{fig:vault-lifecycle}
\end{figure}

Figure~\ref{fig:vault-lifecycle} shows the switch's lifecycle. The decisive
property is that the transition from \emph{Active} to \emph{Triggered} is not an
action anyone takes---it is simply the passage of time making the
\texttt{isDeceased} predicate true. No transaction, no operator, and no
cooperation are required for the switch to fire. This is what makes it
unstoppable.

When the switch triggers, the Lit network---whose access-control condition reads
\texttt{isDeceased(source)}---will release the AES key to a requester, who fetches
the ciphertext from Arweave and decrypts it. The off-chain key release and the
on-chain trigger are bound together by the access-control condition, so that the
key becomes available exactly when, and only when, the source goes silent.

\section{Module B: Identity and Reputation}
\label{sec:design-identity}

Module B answers the credibility problem. Its on-chain component, the
\texttt{WhistleblowerRegistry}, stores, per source, a list of proof
identifiers/hashes, a verified flag, a registration timestamp, and a published
stealth meta-address. By principle P3, the \emph{full} zkTLS proof never goes
on-chain---only its hash (or, in the trustless path, the claim identifier
verified on-chain)---while the proof JSON lives off-chain.

The registry deliberately supports three verification modes of decreasing
strength, so the system can be honest (principle P6) about what a ``verified''
badge means in a given deployment:

\begin{enumerate}
  \item \textbf{On-chain verified (trustless).} The full Reclaim proof is passed
  to a deployed Reclaim verifier contract, which checks the witness signatures
  on-chain. Verification accrues to the proof's \emph{owner}, not the caller, so a
  valid proof cannot be replayed to verify a different account. This is the
  production path and gives the verified flag cryptographic meaning with no
  trusted party.
  \item \textbf{Owner attestation.} A designated verifier (the contract owner,
  e.g.\ a backend or a Lit Action) attests verification after checking a proof
  off-chain---used where no verifier contract is deployed for the chain or
  provider.
  \item \textbf{Trust-on-submit (development only).} The first submitted proof
  hash marks the submitter verified. This is convenient for local demonstration
  and is explicitly \emph{not} a credibility guarantee; it defaults off on live
  networks.
\end{enumerate}

This layered design is itself a contribution: rather than pretending a single
mechanism suits all deployments, the registry makes the trust assumption explicit
and configurable, and the user interface selects the strongest available mode for
the active network.

\section{Module C: The Marketplace}
\label{sec:design-marketplace}

Module C lets sources be compensated. The \texttt{Marketplace} contract holds
listings (a description hash, an Arweave reference to the encrypted payload, a
minimum bid, and a verified flag \emph{read from the registry}, not self-asserted)
and escrowed bids. A journalist places a bid by sending ether, which the contract
holds; the source accepts a bid, at which point the contract deducts a fixed
platform fee and forwards the remainder to a one-time \emph{stealth address}
supplied by the source. Unaccepted bids can be withdrawn for a full refund.

\begin{figure}[htbp]
\centering
\begin{tikzpicture}[
  font=\sffamily\scriptsize,
  actor/.style={draw=zkwaccent, thick, rounded corners, fill=zkwbg, minimum width=2.3cm, minimum height=0.9cm, align=center},
  op/.style={draw=zkwframe, rounded corners, fill=white, minimum width=2.6cm, minimum height=0.85cm, align=center},
  >={Stealth[length=2mm]}
]
\node[actor] (recv) {Source\\(receiver)};
\node[op, right=2.4cm of recv] (meta) {publish\\meta-address};
\node[actor, right=2.4cm of meta] (send) {Buyer\\(sender)};

\node[op, below=1.1cm of meta] (eph) {ECDH with\\ephemeral key};
\node[op, below=1.1cm of eph] (stealth) {one-time\\stealth address};
\node[op, left=2.6cm of stealth] (scan) {scan with\\viewing key};

\draw[->, thick] (recv) -- (meta);
\draw[->, thick] (meta) -- (send);
\draw[->, thick] (send) -- (eph);
\draw[->, thick] (eph) -- (stealth);
\draw[->, thick] (stealth) -- node[below, font=\sffamily\tiny] {payout (minus fee)} (scan);
\draw[->, thick] (scan) -- node[left, font=\sffamily\tiny] {recover funds} (recv);
\draw[->, dashed, zkwcomment] (eph) -| node[near start, above, font=\sffamily\tiny] {announce ephemeral pubkey} (scan);
\end{tikzpicture}
\caption{Stealth-address payment flow (ERC-5564). The buyer derives a fresh
one-time address from the source's published meta-address; only the source, using
their viewing and spending keys, can locate and spend the funds.}
\label{fig:stealth-flow}
\end{figure}

Figure~\ref{fig:stealth-flow} illustrates the payment. Because the payout goes to
a freshly derived address with no on-chain link to the source's published
identity, an observer sees only a transfer to an unfamiliar account. This
satisfies principle P5's requirement that reward not break anonymity. Coupling the
marketplace's verified flag to the registry (rather than to a self-asserted
parameter) ensures a ``Verified Source'' badge always reflects genuine on-chain
reputation.

\section{Data Model: On-Chain versus Off-Chain}
\label{sec:design-datamodel}

The placement of data is a security decision, not merely an efficiency one.
Table~\ref{tab:datamodel} enumerates where each kind of data lives and why,
making the ``no plaintext on-chain'' principle (P3) concrete.

\begin{table}[htbp]
\centering
\caption{What is stored where, and why.}
\label{tab:datamodel}
\small
\begin{tabular}{@{}p{4.4cm}p{3.0cm}p{6.0cm}@{}}
\toprule
\textbf{Data} & \textbf{Location} & \textbf{Rationale} \\
\midrule
File plaintext            & Client only (RAM) & Never persisted anywhere (P2). \\
File ciphertext           & Arweave           & Permanent, censorship-resistant; useless without the key. \\
AES key (plaintext)       & Client only (RAM) & Exported, then immediately sealed by Lit. \\
AES key (sealed)          & Arweave manifest  & Only the Lit network can unseal it, and only under the ACC. \\
Heartbeat \& interval     & EVM (DeadMansSwitch) & Must be public/objective for the ACC to read trustlessly. \\
Arweave / Lit references  & EVM               & Coordination pointers; reveal nothing sensitive. \\
zkTLS proof (full JSON)   & Off-chain         & Large; only its hash/identifier needs to be anchored. \\
Proof hash / identifier   & EVM (Registry)    & Anchors credibility; reveals nothing about identity. \\
Stealth meta-address      & EVM (Registry)    & Public by design; unlinkable to derived payments. \\
Bids \& payment state     & EVM (Marketplace) & Needed for escrow; payout goes to a stealth address. \\
\bottomrule
\end{tabular}
\end{table}

\section{End-to-End Workflows}
\label{sec:design-workflows}

Two workflows tie the modules together: arming a vault and releasing it.

\paragraph{Arming.} The source (i) selects a file, which is encrypted in the
browser with a fresh AES-256-GCM key; (ii) the AES key is sealed by Lit under the
\texttt{isDeceased(source)} access-control condition; (iii) a self-describing
manifest---ciphertext, IV, sealed key, and metadata, but never plaintext---is
uploaded to Arweave; and (iv) the source calls \texttt{createSwitch} with the
chosen interval, the Arweave reference, the Lit reference, and the recipient.
Thereafter they call \texttt{checkIn} within each interval.

\paragraph{Releasing.} Figure~\ref{fig:release-seq} shows the release. When the
interval lapses, \texttt{isDeceased(source)} becomes true. A requester (the named
recipient, or anyone, for a public release) fetches the manifest from Arweave,
authenticates to Lit with their wallet, and asks Lit to unseal the AES key. The
Lit nodes independently evaluate the access-control condition against the chain;
finding it true, they each contribute a decryption share, and the combined key is
returned to the requester, who decrypts the payload locally.

\begin{figure}[htbp]
\centering
\begin{tikzpicture}[
  font=\sffamily\scriptsize,
  party/.style={draw=zkwaccent, thick, fill=zkwbg, rounded corners, minimum width=2.1cm, minimum height=0.8cm, align=center},
  >={Stealth[length=2mm]}
]
% lifelines
\node[party] (req) at (0,0) {Requester};
\node[party] (arw) at (3.7,0) {Arweave};
\node[party] (lit) at (7.4,0) {Lit network};
\node[party] (chain) at (11.0,0) {DeadMansSwitch};
\foreach \n in {req,arw,lit,chain} {
  \draw[zkwframe, thick] (\n.south) -- ++(0,-5.4);
}
% messages
\draw[->, thick] ($(req.south)+(0,-0.5)$) -- node[above, font=\sffamily\tiny] {fetch manifest} ($(arw.south)+(0,-0.5)$);
\draw[<-, thick] ($(req.south)+(0,-1.1)$) -- node[above, font=\sffamily\tiny] {ciphertext + sealed key} ($(arw.south)+(0,-1.1)$);
\draw[->, thick] ($(req.south)+(0,-1.9)$) -- node[above, font=\sffamily\tiny] {decrypt request (SIWE auth)} ($(lit.south)+(0,-1.9)$);
\draw[->, thick] ($(lit.south)+(0,-2.5)$) -- node[above, font=\sffamily\tiny] {read isDeceased(source)} ($(chain.south)+(0,-2.5)$);
\draw[<-, thick] ($(lit.south)+(0,-3.1)$) -- node[above, font=\sffamily\tiny] {true} ($(chain.south)+(0,-3.1)$);
\node[align=center, font=\sffamily\tiny, text=zkwcomment] at ($(lit.south)+(0,-3.7)$) {threshold nodes\\combine shares};
\draw[<-, thick] ($(req.south)+(0,-4.4)$) -- node[above, font=\sffamily\tiny] {unsealed AES key} ($(lit.south)+(0,-4.4)$);
\node[align=center, font=\sffamily\tiny, text=zkwcomment] at ($(req.south)+(0,-5.0)$) {AES-GCM decrypt\\locally};
\end{tikzpicture}
\caption{Release sequence. The Lit network releases the key only after
independently confirming the on-chain trigger condition; decryption itself happens
on the requester's device.}
\label{fig:release-seq}
\end{figure}

This completes the design. The next chapter describes how each element is built.
